\documentclass{article}

\usepackage{booktabs}
\usepackage{geometry}
\usepackage{array}
\usepackage{hyperref}

\geometry{a4paper, margin=2.5cm}

\title{E2 MLP Experimental Results and Analysis}
\author{}
\date{}

\begin{document}

\maketitle

\section{Experimental Setup}

This experiment is based on the MLP multi-label classification task under the
\texttt{21/E2} directory. Three loss functions are compared:
\texttt{BCEWithLogitsLoss}, Focal Loss, and Asymmetric Loss (ASL).
The corresponding experiment folders are
\texttt{mlp\_runs\_bce},
\texttt{mlp\_runs\_focal\_g2\_a025}, and
\texttt{mlp\_runs\_asl\_gn4\_gp1\_c005}.
The analysis mainly uses the final test results from each \texttt{test\_metrics.json} file.

The input representation is formed by concatenating six extracted feature groups:
\texttt{Normalized\_CORR},
\texttt{Normalized\_WT},
\texttt{Normalized\_CH},
\texttt{Normalized\_EDH},
\texttt{Normalized\_CM55}, and
\texttt{BoW\_int}.
The training labels are loaded from \texttt{dataset/database\_labels.npy}, and the test
labels are loaded from \texttt{dataset/test\_labels.npy}.

All three runs share the same MLP architecture and training configuration:

\begin{itemize}
    \item MLP hidden dimensions: \([1024, 512, 256]\)
    \item Activation: GELU
    \item Dropout: 0.3
    \item Batch size: 32
    \item Epochs: 50
    \item Learning rate: \(1 \times 10^{-4}\)
    \item Weight decay: \(1 \times 10^{-5}\)
    \item Validation ratio: 0.2
    \item Random seed: 42
    \item Prediction threshold: 0.5
    \item Top-\(k\): 5
    \item Model selection metric: validation mAP
    \item Standardization: disabled
\end{itemize}

Therefore, this experiment mainly evaluates the effect of different loss functions
under the same MLP configuration.

\section{Test Results}

\begin{table}[htbp]
\centering
\caption{Test performance comparison of three loss functions in E2}
\label{tab:e2_loss_comparison_en}
\begin{tabular}{lrrrrrrrr}
\toprule
Loss & Test Loss & Accuracy & Micro-F1 & Macro-F1 & mAP & Example-F1 & Jaccard & P@5 \\
\midrule
BCE
& 0.2495 & 13.38 & 62.60 & 48.18 & 64.48 & 54.51 & 44.95 & 50.68 \\
Focal
& 0.0254 & 8.19 & 47.17 & 28.80 & 64.66 & 37.12 & 29.60 & 50.69 \\
ASL
& 0.0460 & 10.52 & 65.91 & 59.83 & 64.83 & 62.84 & 50.76 & 50.85 \\
\bottomrule
\end{tabular}
\end{table}

\begin{table}[htbp]
\centering
\caption{Recall@5 comparison}
\label{tab:e2_recall_at_5_en}
\begin{tabular}{lr}
\toprule
Loss & Recall@5 \\
\midrule
BCE & 78.86 \\
Focal & 78.90 \\
ASL & 78.89 \\
\bottomrule
\end{tabular}
\end{table}

Accuracy, F1 scores, mAP, Example-F1, Jaccard, P@5, and Recall@5 are reported as
percentages. Test Loss is kept in its original scale. Since BCE, Focal Loss, and
ASL define and scale their objectives differently, their loss values should not be
directly compared across methods.

\section{Analysis}

In terms of mAP, the three loss functions show only small differences.
ASL achieves the highest mAP at 64.83\%, followed by Focal Loss at 64.66\% and BCE
at 64.48\%. This suggests that all three methods have similar ranking ability for
relevant labels.

However, ASL shows a much clearer advantage after converting probabilities into
binary predictions with the fixed threshold of 0.5. ASL obtains the best Micro-F1,
Macro-F1, Example-F1, and Jaccard scores. Its Micro-F1 reaches 65.91\%, Macro-F1
reaches 59.83\%, Example-F1 reaches 62.84\%, and Jaccard reaches 50.76\%.
These results indicate that ASL produces a better balance between positive and
negative label predictions under the fixed-threshold setting.

BCE is a stable baseline. It reaches 62.60\% Micro-F1, 48.18\% Macro-F1, and
54.51\% Example-F1. However, compared with ASL, BCE is much weaker in Macro-F1,
which suggests that it is less effective for minority or difficult labels. In
multi-label classification, label imbalance is common, and BCE may be dominated by
frequent negative labels because it treats all labels with the same binary
classification objective.

Focal Loss shows a different behavior. Its mAP and P@5 are very close to the other
two methods, and its mAP is slightly higher than BCE. Nevertheless, its Micro-F1,
Macro-F1, Example-F1, and Jaccard are much lower. This means that Focal Loss does
not significantly damage label ranking, but its output scores may be less well
calibrated for the fixed threshold of 0.5. If Focal Loss is further studied,
threshold tuning or class-wise threshold search should be considered.

Recall@5 is almost identical across the three methods, around 78.9\%. This further
shows that the main performance difference does not come from top-\(k\) ranking,
but from the final binary label decision.

\section{Conclusion}

Overall, ASL is the best loss function in this E2 experiment. Although its mAP is
only about 0.35 percentage points higher than BCE, it brings much larger
improvements in F1-based metrics: approximately 3.31 percentage points in
Micro-F1, 11.65 percentage points in Macro-F1, and 8.33 percentage points in
Example-F1. This indicates that ASL is more suitable for the current multi-label
classification task, especially under label imbalance and fixed-threshold
prediction.

BCE should be kept as a reliable baseline for future comparisons. Focal Loss is
not recommended as the final choice under the current configuration, because it
preserves similar mAP and top-\(k\) performance but performs poorly after
thresholding. Future work on Focal Loss should focus on threshold tuning, different
\(\gamma\) and \(\alpha\) settings, and validation-based class-wise threshold search.

\end{document}
